% ============================================================================
% paper_body.tex — everything after the Introduction.
%   §2  Constructing objective-specific demand maps   <- WRITTEN (uses Fig. 1)
%   §3  A taxonomy of demand                           <- SKETCH (Table 2)
%   §4  Different demand, different constellations      <- SKETCH (Figs. 2-3)
%   §5  Toward first-class demand modeling (agenda)     <- SKETCH
%   §6  Related work                                    <- SKETCH
%   §7  Conclusion                                      <- SKETCH
%
% Numbers below are from the live pipeline today (1,500 metros >=300k pop,
% 72x144 equal-area grid). population/gdp/road are REAL fetched data.
% nightlight is the PROXY (no VIIRS raster supplied) — see the boxed caveat
% in §2. Recompute nightlight-dependent numbers once the real raster is in.
%
% Figure asset: fig1_demand_signals.pdf (place in your figs/ dir).
% ============================================================================

\section{Constructing objective-specific demand maps}
\label{sec:signals}

If demand is part of the specification, then the first practical question is
whether objective-specific demand maps are actually \emph{different}
maps---or whether, once gridded and normalized, they collapse into
lightly-reweighted versions of the same population surface. This section
builds four such maps from independent public signals and shows they are
not.

\paragraph{A common substrate.}
All four maps share one spatial substrate so that any difference between
them comes from the signal, not the tessellation. We bin the Earth into an
equal-area grid of $n_{\text{lat}}\times n_{\text{lon}}$ cells, uniform in
$\sin(\text{latitude})$ so that a degree of longitude near the pole does not
over-count relative to one at the equator.\footnote{A uniform lat/lon grid
over-weights high-latitude cells; binning uniformly in $\sin(\text{lat})$
carves cells of approximately equal physical area. This is an
implementation detail, not a contribution, but it matters: an unequal-area
grid silently biases every demand centroid poleward.} Onto this grid we
place \todo{$N$; today 1{,}500} metros of population $\ge$\,\todo{300k} from
GeoNames, each carrying population, per-capita GDP (World Bank), a nighttime
radiance sample, and a road-network-reach value. Each demand model is then a
function from this metro table to a per-cell weight; every value in the
pipeline is tagged with its provenance (fetched vs.\ proxied) so a reader
can see exactly which numbers are real signal and which are placeholders.

\paragraph{Four signals, four objectives.}
Figure~\ref{fig:signals} shows the four maps. Each targets a different
objective, and---crucially---each is built from a signal that is
\emph{independent} of the others rather than a transform of population:

\begin{itemize}[leftmargin=1.2em,itemsep=1pt]
  \item \textbf{Population} (universal broadband): weight $\propto$ people.
        The baseline every other map is measured against.
  \item \textbf{GDP} (economically weighted capacity): weight $\propto$
        population $\times$ per-capita GDP, a proxy for traffic that tracks
        spending power. Appropriate for backbone backup and inter-metro
        long-haul, where the customers are the world's economic centers.
  \item \textbf{Nightlight gap} (underserved / under-electrified): weight
        $\propto$ population $\times\,(1-\text{rank}(\text{radiance}))$, so
        weight moves to places that are populated but \emph{dark} at
        night---a proxy for regions terrestrial infrastructure has
        under-served.
  \item \textbf{Road gap} (rural / terrestrial-access gap): weight
        $\propto$ population $\times\,(1-\text{rank}(\text{road km within
        60\,km}))$, so weight moves to populated places with little nearby
        road network. Roads correlate with terrestrial ISP build-out, so
        ``people, but no roads'' is a direct signal that satellite is the
        only realistic option.
\end{itemize}

The last two maps share a template worth stating explicitly:
\emph{population minus terrestrial-infrastructure reach}, instantiated with
two \emph{independent} proxies for that reach (nighttime light and road
density). This is the map several LSN papers gesture at when they invoke the
``millions of unconnected users'' but do not build; here it is concrete.

\paragraph{The maps genuinely diverge.}
Table~\ref{tab:divergence} quantifies the pairwise disagreement. Even though
all four maps are restricted to the \emph{same} metros---so these numbers
are a conservative floor; adding genuinely rural, sub-metro demand points
would only widen the gap---the objective-specific maps are far from
interchangeable. The GDP map and the road-gap map, the two extremes, have a
Spearman rank correlation of only \todo{0.28}, their 100 highest-demand
cells overlap by a Jaccard index of just \todo{0.21}, and their
demand-weighted centroids sit roughly \todo{2{,}900\,km} apart---one pulled
toward wealthy northern metros, the other toward populous but road-poor
regions nearer the equator. The maps also differ in \emph{shape}: the GDP
map concentrates \todo{20}\% of its total weight in its top 1\% of cells,
while the two gap maps spread weight far more evenly
(\todo{8--12}\%)---exactly the behavior you would want from a map meant to
represent diffuse, underserved demand rather than a handful of megacities.

\begin{figure*}[t]
  \centering
  \includegraphics[width=\textwidth]{figs/fig1_demand_signals.pdf}
  \caption{Four objective-specific demand maps on a common equal-area grid,
  built from independent public signals. Marker size and color encode
  normalized per-cell demand; the red $\times$ is each map's demand-weighted
  centroid. The maps place weight in visibly different
  places---economic-center-heavy (GDP) versus dispersed and
  equatorward (road gap)---despite sharing an identical set of metros.
  \todo{Panel (c) uses a population$\times$wealth \emph{proxy} for
  nightlight; replace with the real VIIRS \texttt{avg\_rade9h} raster before
  submission (see text).}}
  \label{fig:signals}
\end{figure*}

\begin{table}[t]
  \centering\small
  \caption{Pairwise divergence across the four maps (union of nonzero
  cells). Low correlation, low top-100 overlap, and large centroid
  separation all say the same thing: these are different maps.
  \todo{Regenerate the nightlight rows with the real VIIRS raster.}}
  \label{tab:divergence}
  \begin{tabular}{lccc}
    \toprule
    map pair & Spearman $\rho$ & top-100 Jaccard & centroid dist. \\
    \midrule
    population -- gdp            & 0.57 & 0.40 & 1{,}767\,km \\
    population -- road gap       & 0.82 & 0.52 & 1{,}187\,km \\
    gdp -- road gap              & \textbf{0.28} & \textbf{0.21} & \textbf{2{,}954\,km} \\
    \addlinespace
    population -- nightlight\,$\dagger$ & 0.96 & 0.68 & 868\,km \\
    gdp -- nightlight\,$\dagger$        & 0.36 & 0.30 & 2{,}611\,km \\
    road gap -- nightlight\,$\dagger$   & 0.85 & 0.49 & 404\,km \\
    \bottomrule
  \end{tabular}
  \\[2pt]
  {\footnotesize $\dagger$ proxy nightlight; provisional (see caveat).}
\end{table}

% !!!!!!!!!!!!!!!!!!!!!!!!!!!!!!!!!!!!!!!!!!!!!!!!!!!!!!!!!!!!!!!!!!!!!!!!!!!!!!
% CAVEAT — READ BEFORE SUBMITTING. The nightlight column in this run is the
% pipeline's population x wealth PROXY, because no VIIRS raster was supplied.
% Consequence: nightlight correlates 0.96 with population *by construction*,
% which UNDERSELLS the thesis (it makes one of your four maps look redundant).
% Do NOT feature panel (c) or quote any nightlight number until you drop in
% the real raster:
%     python make_demand_map.py --model nightlight_gap \
%         --nightlight-raster <...>.avg_rade9h.tif --out fig1c.png
% With real radiance, nightlight_gap should DIVERGE from population (dark-but-
% populated regions up-weighted). The gdp--road_gap row is proxy-free and is
% your safe headline tonight. Upload the .tif here and I'll regenerate Fig.1.
% !!!!!!!!!!!!!!!!!!!!!!!!!!!!!!!!!!!!!!!!!!!!!!!!!!!!!!!!!!!!!!!!!!!!!!!!!!!!!!


% ==========================================================================
\section{A taxonomy of demand}       % ===== SKETCH =====
\label{sec:taxonomy}
% GOAL (~0.75 pp): make the scenario->demand mapping from the intro
% systematic, and show each objective needs its own signal(s) and, sometimes,
% its own *representation* (point vs. pair/transit).

% Opening move: a demand model has three parts a paper should have to state
% — WHERE (spatial support), HOW MUCH (weighting), and WHEN (temporal
% structure). Most published models declare only the first two, implicitly.

\begin{table}[t]
  \centering\small
  \caption{Objectives induce demand models. Each row is a distinct notion of
  demand with distinct data sources; the ``rep.'' column flags whether the
  objective needs point demand or pair/transit demand — a representational
  fork most single-map pipelines cannot express. \todo{fill/verify}}
  \label{tab:taxonomy}
  \begin{tabular}{llll}
    \toprule
    objective & signal / proxy & data source & rep. \\
    \midrule
    universal broadband   & population            & GeoNames/WorldPop      & point \\
    econ.\ capacity       & pop $\times$ GDPpc    & World Bank             & point \\
    rural / underserved   & pop $-$ infra reach   & VIIRS, roads (OSM)     & point \\
    backbone augment.     & cable topology        & TeleGeography          & pair  \\
    low-latency long-haul & metro--metro geometry & pop centers + gt-circle & pair \\
    disaster response     & outage history        & IODA, news, \todo{}    & point(t) \\
    maritime / aviation   & lane / corridor density & AIS, ADS-B           & point(t) \\
    \bottomrule
  \end{tabular}
\end{table}

% Body (2-3 short paras):
%  - Walk 2-3 rows to show the induced model really differs, using Fig. 1 as
%    the worked example for the point-demand rows (broadband/econ/rural).
%  - The representational fork: backbone & long-haul are PAIR demand. Point
%    maps can't say "traffic FROM region A TO region B." Flag this as the
%    honest boundary of what Fig.1-style maps do, and forward-ref §5. (This
%    is also the clean answer to Ram's "econ vs long-haul?" question: econ =
%    where capacity terminates (point); long-haul = which pairs it connects.)
%  - Temporal column: disaster/maritime/aviation carry a WHEN that none of
%    our four static maps encode. Name it as a gap, not a solved thing.
% Punchline sentence: "No single map serves two rows of this table; that is
% the taxonomy's whole content."


% ==========================================================================
\section{Different demand, different constellations}   % ===== SKETCH =====
\label{sec:crosseval}
% GOAL (~1.5 pp): the empirical payoff. Same synthesizer, only the demand map
% changes -> constellations don't transfer. This is the "money" section.

% --- Setup paragraph (state the controlled-variable discipline up front) ---
% Same framework for every run: matching-pursuit selection of repeat-ground-
% track orbits (TinyLEO), SAME satellite budget N, SAME altitude limits, SAME
% launch/inclination constraints. The ONLY thing that changes between runs is
% the demand map fed in. No SS-orbits (deliberately — avoids a reviewer
% detour into whether SS-orbits are valid; cite TinyLEO's RGT machinery and
% move on). Normalize every map to equal total demand and synthesize to a
% fixed budget N, reporting % demand satisfied — otherwise satellite counts
% aren't comparable (this is the "balancing/weighting" fix).

\begin{figure}[t]\centering
  % \includegraphics[width=\columnwidth]{figs/fig2_crosseval.pdf}
  \fbox{\parbox{0.9\columnwidth}{\centering\vspace{3em}
  \textbf{Fig().~2 (TODO): cross-evaluation heatmap.}\\[2pt]
  rows = demand model synthesized for; cols = demand model evaluated
  against; cell = \% demand met at fixed budget $N$. Diagonal should
  dominate.\vspace{3em}}}
  \caption{A constellation serves the demand it was built for and
  under-serves the rest. \todo{generate from synthesis runs}}
  \label{fig:crosseval}
\end{figure}

% --- Result paragraph ---
% Read Fig. 2: strong diagonal, off-diagonal as low as \todo{Y}%. State the
% single headline number (worst off-diagonal transfer) in the first sentence.

% --- Orbital-structure paragraph + Fig. 3 (inclination/RAAN per objective) --
% Per Ram: the orbital differences are EVIDENCE for the thesis, not the
% contribution. Explain physically: rural/road-gap spreads planes across
% low-latitude land masses; GDP concentrates over northern mid-latitudes;
% etc. Report inclination, altitude, sats/plane per run in a small table.

% --- Guardrail paragraph (answer Chris's "of course it specializes") ---
% Pre-empt head-on: yes, tailored designs are supposed to fit their targets;
% specialization working is the MECHANISM, not a rebuttal. Because it works:
% (a) an evaluation claiming objective O but optimizing demand D!=O has not
% established what it claims; (b) reported comparisons are demand-conditional.
% [If the ranking-flip experiment lands, its figure goes here as the sharpest
% form of (b): the "best" design inverts when the demand map changes.]


% ==========================================================================
\section{Toward first-class demand modeling}          % ===== SKETCH =====
\label{sec:agenda}
% GOAL (~0.75 pp): the perspective-paper payload — open problems. Keep each
% to 2-3 sentences; this is an agenda, not a solution.

% (1) Transit demand without artificial nodes. Backbone/long-haul are pair
%     demand, but LSN users aren't nodes. How to define endpoints for a
%     gravity-style pair model over a diffuse client population? (Directly
%     picks up the §3 representational fork.)
% (2) Temporal structure. Sector-dependent diurnal patterns (education vs
%     agriculture vs transport); disaster as transient spikes. Demand as a
%     spatiotemporal field, not a static map. (Ties to the "and when" in the
%     intro's demand-model definition.)
% (3) Validation without ground truth. The populations that motivate LSNs
%     generate no measurable traffic today, so maps can't be checked against
%     observation. What does falsification even look like? (cross-proxy
%     agreement; natural experiments when new coverage arrives; sensitivity.)
% (4) Fusion & provenance. Corporate filings, day-side RF observation, AIS,
%     ADS-B, cable-landing topology, outage histories — combined WITH the
%     provenance discipline from §2 so a map's real-vs-proxied composition is
%     legible.
% (5) A benchmark. A public suite of objective-LABELED demand maps, so papers
%     report per-objective results against declared demand instead of one
%     unlabeled number. (This operationalizes the whole thesis; our four maps
%     + pipeline are a first seed, released with the paper.)


% ==========================================================================
\section{Related work}                                % ===== SKETCH =====
\label{sec:related}
% ~0.4 pp, three clusters:
%  - LSN design/synthesis: TinyLEO, MegaReduce, SS-Plane; simulators
%    (Hypatia, StarPerf-style). How each represents demand -> our critique.
%  - Terrestrial demand modeling: traffic matrices, gravity models — and why
%    they don't port (no nodes; the unconnected have no "before" traffic).
%  - Connectivity / digital-divide mapping: nightlights literature, Ookla,
%    ITU/GSMA coverage. We reuse these as SIGNALS but for a new purpose.


% ==========================================================================
\section{Conclusion}                                  % ===== SKETCH =====
\label{sec:conclusion}
% 2-3 sentences. Restate: demand is part of the specification; an "optimal"
% constellation without a declared demand model is ill-posed. The ask:
% declare your demand model the way you declare topology and workload. Point
% at the released maps + benchmark seed as the way to start doing so.
